Can graph learning over a road network improve on independent forecasts of each road
segment? We study 5--60-minute speed forecasting at 207 sensors of the METR-LA
benchmark under a strict protocol: chronological splits, the same one-hour history
for every model, masked metrics at every horizon, and peak, off-peak and congested
regimes. A strong per-segment baseline of 207 multi-output XGBoost models reaches a
mean absolute error (MAE) of 3.20\,mph. On out-of-fold training residuals, we show
that its errors are spatially correlated and predictable from the state of
neighbouring segments at forecast time, which is the information an independent
model cannot use. We then propose ST-GAT, a spatio-temporal graph attention network
over the fixed road graph with separate downstream and upstream attention,
outage masks and cyclic time inputs, and a hybrid in which graph attention learns a
residual correction to the XGBoost forecasts, trained on leakage-free cross-fitted
base predictions. ST-GAT lowers the MAE to 3.04\,mph ($-5\,\%$; 3.47\,mph at
60\,min), and the hybrid to 3.06\,mph, with the best 15-minute accuracy
(2.72\,mph).
Both gains over the independent models are significant at every horizon
(Diebold--Mariano, $p<0.001$) and largest in peak periods and congestion. Ablations
isolate the value of the graph, of attention, of the road prior, of outage masking,
of the residual connection and of cross-fitting. An offline, simulated signal-plan
experiment shows that network-aware forecasts let a controller anticipate 1.5--1.8
times as many congestion onsets as the independent model; we make no claim of live
control.
